\documentclass[12pt,a4paper]{article}
\usepackage[utf8]{inputenc}
\usepackage{amsmath}
\usepackage{amssymb}
\usepackage{geometry}
\usepackage{array}
\usepackage{enumitem}
\usepackage{tikz}
\geometry{left=2cm,right=2cm,top=2cm,bottom=2cm}

\title{\textbf{Grade 8 Mathematics - Baseline Test\\Marking Memorandum}}
\author{CAPS Curriculum 2026}
\date{}

\begin{document}

\maketitle

\vspace{0.5cm}
\noindent\textbf{Total Marks: 20}

\vspace{0.3cm}
\hrule
\vspace{0.5cm}

\section*{Section A: Grade 7 Recap [8 marks]}

\subsection*{Question 1: Integers and Operations [3 marks]}

\textbf{(a) Calculate: $-7 + 15 - (-4)$ \hfill [2 marks]}

\vspace{0.2cm}
\noindent\textbf{Solution:}

$-7 + 15 - (-4)$ \checkmark \quad \textit{(rewrite subtracting negative as addition)}

$= -7 + 15 + 4$ \checkmark \quad \textit{(correct simplification)}

$= 12$ \checkmark\checkmark \quad \textit{(final answer)}

\vspace{0.3cm}
\noindent\textbf{Mark allocation:}
\begin{itemize}
    \item 1 mark: Correct method (rewriting $-(-4)$ as $+4$)
    \item 1 mark: Correct final answer
\end{itemize}

\vspace{0.4cm}

\textbf{(b) Calculate: $(-3) \times 4 \times (-2)$ \hfill [1 mark]}

\vspace{0.2cm}
\noindent\textbf{Solution:}

$(-3) \times 4 \times (-2)$

$= -12 \times (-2)$ \checkmark \quad \textit{(first multiplication)}

$= 24$ \checkmark \quad \textit{(final answer, negative $\times$ negative = positive)}

\vspace{0.3cm}
\noindent\textbf{Mark allocation:}
\begin{itemize}
    \item 1 mark: Correct final answer
\end{itemize}

\subsection*{Question 2: Fractions and Percentages [3 marks]}

\textbf{(a) Calculate: $\dfrac{2}{3} + \dfrac{3}{4}$ \hfill [2 marks]}

\vspace{0.2cm}
\noindent\textbf{Solution:}

$\dfrac{2}{3} + \dfrac{3}{4}$

$= \dfrac{2 \times 4}{3 \times 4} + \dfrac{3 \times 3}{4 \times 3}$ \checkmark \quad \textit{(find LCD = 12)}

$= \dfrac{8}{12} + \dfrac{9}{12}$ \checkmark \quad \textit{(correct conversion)}

$= \dfrac{17}{12}$ \checkmark\checkmark \quad \textit{(final answer)}

$= 1\dfrac{5}{12}$ \quad \textit{(alternative form)}

\vspace{0.3cm}
\noindent\textbf{Mark allocation:}
\begin{itemize}
    \item 1 mark: Finding common denominator (LCD = 12)
    \item 1 mark: Correct final answer
\end{itemize}

\vspace{0.4cm}

\textbf{(b) Convert $0.35$ to a percentage and simplify $\dfrac{7}{20}$ to a decimal. \hfill [1 mark]}

\vspace{0.2cm}
\noindent\textbf{Solution:}

$0.35 = 35\%$ \checkmark \quad \textit{(multiply by 100)}

$\dfrac{7}{20} = 0.35$ \checkmark \quad \textit{(divide 7 by 20)}

\vspace{0.3cm}
\noindent\textbf{Mark allocation:}
\begin{itemize}
    \item 1 mark: Both conversions correct
\end{itemize}

\subsection*{Question 3: Basic Algebra [2 marks]}

\textbf{(a) Simplify: $5x + 3y - 2x + 7y$ \hfill [1 mark]}

\vspace{0.2cm}
\noindent\textbf{Solution:}

$5x + 3y - 2x + 7y$

$= (5x - 2x) + (3y + 7y)$ \checkmark \quad \textit{(group like terms)}

$= 3x + 10y$ \checkmark \quad \textit{(final answer)}

\vspace{0.3cm}
\noindent\textbf{Mark allocation:}
\begin{itemize}
    \item 1 mark: Correct simplification
\end{itemize}

\vspace{0.4cm}

\textbf{(b) Solve for $x$: $3x + 7 = 22$ \hfill [1 mark]}

\vspace{0.2cm}
\noindent\textbf{Solution:}

$3x + 7 = 22$

$3x = 22 - 7$ \checkmark \quad \textit{(subtract 7 from both sides)}

$3x = 15$

$x = \dfrac{15}{3}$ \checkmark \quad \textit{(divide both sides by 3)}

$x = 5$

\vspace{0.3cm}
\noindent\textbf{Mark allocation:}
\begin{itemize}
    \item 1 mark: Correct solution
\end{itemize}

\newpage
\section*{Section B: Grade 8 Content (Up to August 2026) [12 marks]}

\subsection*{Question 4: Exponents and Powers [3 marks]}

\textbf{(a) Simplify: $2^3 \times 2^4$ \hfill [1 mark]}

\vspace{0.2cm}
\noindent\textbf{Solution:}

$2^3 \times 2^4 = 2^{3+4}$ \checkmark \quad \textit{(add exponents when multiplying same base)}

$= 2^7$

\vspace{0.3cm}
\noindent\textbf{Mark allocation:}
\begin{itemize}
    \item 1 mark: Correct application of exponent law
\end{itemize}

\vspace{0.4cm}

\textbf{(b) Simplify: $\dfrac{x^8}{x^3}$ \hfill [1 mark]}

\vspace{0.2cm}
\noindent\textbf{Solution:}

$\dfrac{x^8}{x^3} = x^{8-3}$ \checkmark \quad \textit{(subtract exponents when dividing same base)}

$= x^5$

\vspace{0.3cm}
\noindent\textbf{Mark allocation:}
\begin{itemize}
    \item 1 mark: Correct application of exponent law
\end{itemize}

\vspace{0.4cm}

\textbf{(c) Write $3\,500\,000$ in scientific notation. \hfill [1 mark]}

\vspace{0.2cm}
\noindent\textbf{Solution:}

$3\,500\,000 = 3.5 \times 10^6$ \checkmark

\vspace{0.3cm}
\noindent\textbf{Mark allocation:}
\begin{itemize}
    \item 1 mark: Correct scientific notation
\end{itemize}

\subsection*{Question 5: Algebraic Expressions and Equations [3 marks]}

\textbf{(a) Expand: $3(2x - 5) + 4(x + 2)$ \hfill [2 marks]}

\vspace{0.2cm}
\noindent\textbf{Solution:}

$3(2x - 5) + 4(x + 2)$

$= 6x - 15 + 4x + 8$ \checkmark\checkmark \quad \textit{(correct expansion of both brackets)}

$= (6x + 4x) + (-15 + 8)$ \quad \textit{(group like terms)}

$= 10x - 7$ \checkmark \quad \textit{(final simplified answer)}

\vspace{0.3cm}
\noindent\textbf{Mark allocation:}
\begin{itemize}
    \item 1 mark: Correct expansion of brackets
    \item 1 mark: Correct simplification
\end{itemize}

\vspace{0.4cm}

\textbf{(b) Solve for $x$: $2(x + 3) = 14$ \hfill [1 mark]}

\vspace{0.2cm}
\noindent\textbf{Solution:}

$2(x + 3) = 14$

$x + 3 = \dfrac{14}{2}$ \checkmark \quad \textit{(divide both sides by 2)}

$x + 3 = 7$

$x = 7 - 3$

$x = 4$ \checkmark \quad \textit{(final answer)}

\vspace{0.3cm}
\noindent\textbf{Mark allocation:}
\begin{itemize}
    \item 1 mark: Correct solution
\end{itemize}

\subsection*{Question 6: Geometry - Pythagoras and Area [3 marks]}

\textbf{(a) A right-angled triangle has sides $a = 6$ cm and $b = 8$ cm. Use Pythagoras to calculate the hypotenuse $c$. \hfill [2 marks]}

\vspace{0.2cm}
\noindent\textbf{Solution:}

According to Pythagoras' theorem:
$c^2 = a^2 + b^2$

$c^2 = 6^2 + 8^2$ \checkmark \quad \textit{(substitution)}

$c^2 = 36 + 64$

$c^2 = 100$

$c = \sqrt{100}$

$c = 10$ cm \checkmark \quad \textit{(final answer with units)}

\vspace{0.3cm}
\noindent\textbf{Mark allocation:}
\begin{itemize}
    \item 1 mark: Correct substitution into formula
    \item 1 mark: Correct final answer with units
\end{itemize}

\vspace{0.4cm}

\textbf{(b) Calculate the area of a circle with radius $r = 7$ cm. (Use $\pi = \dfrac{22}{7}$) \hfill [1 mark]}

\vspace{0.2cm}
\noindent\textbf{Solution:}

$A = \pi r^2$

$A = \dfrac{22}{7} \times 7^2$ \checkmark \quad \textit{(substitution)}

$A = \dfrac{22}{7} \times 49$

$A = 22 \times 7$

$A = 154$ cm$^2$ \quad \textit{(final answer)}

\vspace{0.3cm}
\noindent\textbf{Mark allocation:}
\begin{itemize}
    \item 1 mark: Correct calculation
\end{itemize}

\subsection*{Question 7: Number Patterns [3 marks]}

\noindent Consider the number pattern: $5, 9, 13, 17, \ldots$

\vspace{0.3cm}

\textbf{(a) Write down the common difference. \hfill [1 mark]}

\vspace{0.2cm}
\noindent\textbf{Solution:}

$d = 9 - 5 = 4$ \checkmark

\vspace{0.3cm}
\noindent\textbf{Mark allocation:}
\begin{itemize}
    \item 1 mark: Correct common difference
\end{itemize}

\vspace{0.4cm}

\textbf{(b) Calculate the 10th term. \hfill [1 mark]}

\vspace{0.2cm}
\noindent\textbf{Solution:}

$T_n = a + (n-1)d$

$T_{10} = 5 + (10-1)(4)$ \checkmark \quad \textit{(substitution)}

$T_{10} = 5 + 9 \times 4$

$T_{10} = 5 + 36$

$T_{10} = 41$

\vspace{0.3cm}
\noindent\textbf{Mark allocation:}
\begin{itemize}
    \item 1 mark: Correct 10th term
\end{itemize}

\vspace{0.4cm}

\textbf{(c) Write down the general term (rule) for this pattern. \hfill [1 mark]}

\vspace{0.2cm}
\noindent\textbf{Solution:}

$T_n = a + (n-1)d$

$T_n = 5 + (n-1)(4)$

$T_n = 5 + 4n - 4$

$T_n = 4n + 1$ \checkmark

\vspace{0.3cm}
\noindent\textbf{Mark allocation:}
\begin{itemize}
    \item 1 mark: Correct general term
\end{itemize}

\vspace{1cm}
\hrule
\vspace{0.3cm}
\noindent\textbf{GRAND TOTAL: 20 marks}

\end{document}
